problem:  
Let $\Sigma^n$ ($n\ge7$) be a smooth homotopy sphere (possibly exotic), and consider the set
\[
\mathcal{M}_\lambda(\Sigma^n)=\{\, g\text{ Riemannian metric on }\Sigma^n\mid \sec_g>0,\ \tfrac{K_{\min}(g)}{K_{\max}(g)}\ge \lambda \,\}.
\]
The Differentiable Sphere Theorem (Brendle–Schoen) shows that if $\lambda > \frac14$, then any $g\in\mathcal{M}_\lambda(\Sigma^n)$ implies $\Sigma^n$ is diffeomorphic to the standard sphere $S^n$.  
However, nothing is known for smaller $\lambda$—even the existence of $\sec>0$ metrics on exotic spheres remains open.

**Question:**  
Is there a universal pinching threshold $\lambda^*_n\in(0,\frac14)$, depending only on the dimension, such that if an exotic sphere $\Sigma^n$ admits a positively curved metric with $\frac{K_{\min}}{K_{\max}}>\lambda^*_n$, then $\Sigma^n$ must be diffeomorphic to $S^n$?  

Equivalently, can exotic smooth structures on spheres persist arbitrarily close to the geometric roundness limit, or is there a positive **gap in curvature pinching** separating the regime of possible exotic spheres from that of the standard sphere?  

Moreover, if such a $\lambda^*_n$ exists, can it be characterized dynamically—for example, as the smallest ratio invariant under Ricci flow that separates convergence to a round metric from finite‑time singularity formation on homotopy spheres?

Why is it a "good" problem:  
This question sits squarely at the confluence of **global differential topology**, **Riemannian geometry**, and **geometric analysis**.  
It isolates a precise, quantitative threshold—$\lambda^*_n$—whose value would delineate the analytic boundary between standard and exotic smooth structures among positively curved spheres.  

- It extends the reasoning of the Differentiable Sphere Theorem into a more subtle, as‑yet‑unexplored regime: how round can a positively curved exotic sphere be before topology forces standardness?  
- It ties together two deep mechanisms: the **curvature pinching gap** (a geometric concept) and the **smooth structure of spheres** (a topological concept), inviting new methods from Ricci flow, stability analysis, and curvature pinching evolution.  
- The statement is simple, elegant, and purely geometric: “How close to round can an exotic sphere be?”—yet its resolution would deeply impact the understanding of curvature rigidity and the geometric detection of smooth structure.  

A solution, even partial, would yield a new analytic invariant separating standard from exotic spheres and could bridge differentiable topology with curvature dynamics in a tangible, quantitative way—embodying the clarity, depth, and cross‑field integration that define a truly “good” research‑level problem.